\chapter{Considerações Finais} \label{cap:consideracoes}

\section{Conclusões do Projeto de Formatura}
\label{sec:conclusoes}

% Capítulo exclusivo da monografia final: a parcial não o inclui.
\pendente{Balanço final do trabalho: resultados atingidos e não atingidos, com
justificativas, à luz das quatro perguntas de pesquisa da
Seção~\ref{sec:avaliacao}.}

\section{Contribuições}
\label{sec:contribuicoes}

O trabalho contribui em quatro pontos.

O primeiro é a leitura crítica do gabarito do Racebench, com uma contagem
explícita por unidade, uma errata em que cada correção traz sua evidência e uma
regra de casamento publicada. Sem isso, nenhuma taxa de \textit{recall} dessa literatura
é comparável com outra.

O segundo é a lista explícita de premissas sob as quais a ausência de falsos
negativos vale. As quatro ferramentas revisadas afirmam essa ausência sem
declarar as condições, o que torna a afirmação impossível de verificar
externamente.

O terceiro é a unificação, num só \textit{front end}, das duas formas pelas
quais a literatura enuncia o defeito, hoje tratadas por ferramentas distintas,
com a enumeração de acessos compartilhada e os filtros separados.

O quarto é a aplicação de triagem e reparo assistidos por modelo de linguagem a
defeitos de concorrência, com a ablação das técnicas de \textit{prompt} no
formato de \citeonline{llift}. É a contribuição de maior risco e também a que
não tem precedente em nenhuma das duas literaturas revisadas.

\section{Perspectivas de Continuidade}
\label{sec:continuidade}

Quatro caminhos se abrem a partir do que está especificado aqui.

O primeiro é a inferência automática do modelo de interrupções. O arquivo de
configuração é hoje escrito à mão e é a maior ameaça declarada ao \textit{recall}, já
que um tratador não declarado é invisível para a análise. Inferir os pontos de
entrada a partir da API de registro de interrupções da plataforma eliminaria
essa dependência.

O segundo é o tratamento de reentrância. A premissa de que uma \isr\ não
interrompe a si mesma é a única premissa não conservadora da lista, e removê-la
exige decidir como limitar a profundidade de aninhamento sem perder defeitos.

O terceiro é a avaliação sobre código de produção com vários arquivos. Os dois
\textit{benchmarks} disponíveis são de arquivo único, e a capacidade de ligar
vários módulos existe sem que haja gabarito para exercitá-la.

O quarto é estender a mesma arquitetura a outras classes de defeito de
concorrência, como as que envolvem mais de duas variáveis compartilhadas, que
nenhuma das ferramentas revisadas trata.
